\par\Needspace{6\baselineskip}\medskip
\hypertarget{native-120-article089}{}\pdfbookmark[2]{89.}{article-089}
\begin{center}89.\end{center}
Wir haben nun zwei Aufgaben. Zunächst wollen wir die Unbekannte $g$ möglichst bequem aus der transzendenten Gleichung 12 ermitteln, da diese keine direkte Lösung zulässt. Anschließend wollen wir aus dem so gefundenen Winkel $g$ die Elemente selbst ableiten. Zuvor betrachten wir eine Umformung, \SourceLeaf{121}die die Berechnung der Hilfsgröße $l$ oder $L$ erleichtert und außerdem mehrere später herzuleitende Formeln auf eine elegantere Form bringt.

Dazu führen wir den Hilfswinkel $\omega$ ein, der durch die Formel
\[
\BareRoot[4]{\dfrac{r'}{r}}=\operatorname{tang}(45^0+\omega)
\]
bestimmt wird. Damit ergibt sich
\[
\begin{aligned}
\BareRoot{\dfrac{r'}{r}}+\BareRoot{\dfrac{r}{r'}}
&=2+\{\operatorname{tang}(45^0+\omega)-\operatorname{cotang}(45^0+\omega)\}^2\\[4pt]
&=2+4\operatorname{tang}2\omega^2;
\end{aligned}
\]
und daraus folgt
\[
\begin{aligned}
l&=\frac{\sin\tfrac12f^2}{\cos f}+\frac{\operatorname{tang}2\omega^2}{\cos f}\mathpunct{\raisebox{.45ex}{,}}\\[6pt]
L&=-\frac{\sin\tfrac12f^2}{\cos f}-\frac{\operatorname{tang}2\omega^2}{\cos f}\mathpunct{\cdot}
\end{aligned}
\]
